
\chapter{M002 --- Primitive Surface Construction}

\section{Stage 4 --- Comparison (Phase 4)}

\subsection{Mathematical Proposition}

Certified Gauss reduction replaces an in-plane basis by a gauge-equivalent
canonical representative while preserving the primitive surface lattice
and its certification.

\subsection{Implementation}

The implementation performs Gauss reduction only after constructing the
primitive kernel/B\'ezout certificate and documents preservation of that
certificate as an explicit requirement.

\subsection{Comparison}

\begin{center}
\begin{tabular}{|p{0.44\linewidth}|p{0.34\linewidth}|}
\hline
Mathematical requirement & Implementation status\\
\hline
Primitive lattice preserved & Documented\\
Kernel certificate preserved & Documented\\
Canonical representative selected & Implemented\\
Scientific object unchanged & Consistent with specification\\
\hline
\end{tabular}
\end{center}

\subsection{Evidence}

Evidence reviewed:

\begin{itemize}
\item Stage 2 mathematical specification.
\item Stage 3 implementation review.
\item Implementation algorithm description.
\end{itemize}

\subsection{Audit Result}

The implementation is consistent with the interpretation of Gauss
reduction as a gauge transformation acting on representatives rather
than on the underlying primitive surface.

\subsection{Remaining Question}

Stage 5 should verify that every elementary reduction step preserves the
primitive-kernel certificate and associated lattice invariants.
